\let\negmedspace\undefined
\let\negthickspace\undefined
\documentclass[journal]{IEEEtran}
\usepackage[a5paper, margin=10mm, onecolumn]{geometry}
%\usepackage{lmodern} % Ensure lmodern is loaded for pdflatex
\usepackage{tfrupee} % Include tfrupee package

\setlength{\headheight}{1cm} % Set the height of the header box
\setlength{\headsep}{0mm}     % Set the distance between the header box and the top of the text

\usepackage{gvv-book}
\usepackage{gvv}
\usepackage{cite}
\usepackage{amsmath,amssymb,amsfonts,amsthm}
\usepackage{algorithmic}
\usepackage{graphicx}
\usepackage{textcomp}
\usepackage{xcolor}
\usepackage{txfonts}
\usepackage{listings}
\usepackage{enumitem}
\usepackage{mathtools}
\usepackage{gensymb}
\usepackage{comment}
\usepackage[breaklinks=true]{hyperref}
\usepackage{tkz-euclide} 
\usepackage{listings}
% \usepackage{gvv}                                        
\def\inputGnumericTable{}                                 
\usepackage[latin1]{inputenc}                                
\usepackage{color}                                            
\usepackage{array}                                            
\usepackage{longtable}                                       
\usepackage{calc}                                             
\usepackage{multirow}                                         
\usepackage{hhline}                                           
\usepackage{ifthen}                                           
\usepackage{lscape}
\begin{document}

\bibliographystyle{IEEEtran}
\vspace{3cm}

\title{4.10.12}
\author{EE25BTECH11011 - Navya}
% \maketitle
% \newpage
% \bigskip
{\let\newpage\relax\maketitle}

\renewcommand{\thefigure}{\theenumi}
\renewcommand{\thetable}{\theenumi}
\setlength{\intextsep}{10pt} % Space between text and floats
\textbf{Question}:\\
 The straight line $5x+4y=0$ passes through the point of intersection of the straight lines $x+2y-10 = 0$ and $2x+y+5 = 0$.
 
\textbf{Solution:}\\
The equations of the given lines are
\begin{align}
    \myvec{1 & 2}\vec{x} &= 10 \label{eq1}\\
    \myvec{2 & 1}\vec{x} &= -5 \label{eq2}
\end{align}

Form a single matrix equation (all three lines):
\begin{align}
&\qquad
\myvec{1 & 2\\[4pt] 2 & 1\\[4pt] 5 & 4}\vec{x}
= \myvec{10\\[4pt]-5\\[4pt]0}. \label{system}
\end{align}

To solve, take the first two rows and solve by inverse:
\begin{align}
\myvec{1 & 2\\[4pt]2 & 1}\vec{x} &= \myvec{10\\[4pt]-5}
\end{align}

\begin{align}
\vec{x} 
&= \myvec{1 & 2\\[4pt]2 & 1}^{-1}\myvec{10\\[4pt]-5}= \frac{1}{-3}\myvec{1 & -2\\[4pt]-2 & 1}\myvec{10\\[4pt]-5} = \myvec{-\tfrac{20}{3}\\[6pt]\tfrac{25}{3}}
\end{align}

Now verify that this vector satisfies the third row of (3),
\begin{align}
\myvec{5 & 4}\vec{x} &= 0
\end{align}

\begin{align}
\myvec{5 & 4}\myvec{-\tfrac{20}{3}\\[4pt]\tfrac{25}{3}}
&= \frac{1}{3}(-100+100)= 0 
\end{align}\\

Hence, The straight line $5x+4y=0$ passes through the point of intersection of the straight lines $x+2y-10 = 0$ and $2x+y+5 = 0$. 

\begin{figure}[h!]
    \centering
    \includegraphics[width=0.8\textwidth]{figs/fig8.png}
    \label{fig:division}
\end{figure}
\end{document}